\documentclass[fragile,handout]{beamer}

\newif\iflongonly

%\documentclass{article}
%\usepackage[envcountsect]{beamerarticle}

% class options [notes], [notes=only], [trans], [handout], [compress], [draft]

\input xy
\xyoption{all}


\theoremstyle{plain}
\newtheorem{thm}{Theorem}[section]
\newtheorem{cor}[thm]{Corollary}
\newtheorem{lem}[thm]{Lemma}

\usepackage{bigcenter}
\usepackage{url}

\mode<article> % only for the article version
{
  \usepackage{fullpage}
  \usepackage{pgf}
  \usepackage{hyperref}
  %\setjobnamebeamerversion{talk00.beamer}
}

\usepackage[latin1]{inputenc}
\usepackage[english]{babel}
\usepackage{times}
\usepackage[T1]{fontenc}

\mode<presentation> {
  %\setbeamercovered{transparent}
  \setbeamertemplate{background canvas}[vertical shading][bottom=white,top=structure.fg!25]
  \usetheme{Warsaw}
  \setbeamersize{text margin left=0.7cm}

  %\setbeamertemplate{headline}{}
  %\setbeamertemplate{footline}{}
  %\usetheme{Dresden}
  %\usetheme{Warsaw}
  %\usetheme{Berlin}
  %\useoutertheme{default}
  %\useinnertheme{default}
  %\usefonttheme[onlysmall]{structurebold}
}

%%\href{http://www.famille-blum.org/~william/mscthesis/}

%%%%%%%%%%%%%%%%%%%%% BEGIN OF MY DEFINITIONS
\def\union{\cup}
\def\st{\ |\ }
\def\zor{\vee}
\def\zand{\wedge}
\def\natset{\mathbb{N}}
\def\boolset{\mathbb{B}}
\def\iff{\Longleftrightarrow}
\def\imp{\Rightarrow}
\def\error{\textcolor{red}{\oslash}}
\def\evaluates{\textcolor{orange}{\Downarrow}}
\newcommand{\lsem}{[\![}
\newcommand{\rsem}{]\!]}
\def\sem#1{{\lsem #1 \rsem}}
\def\vec#1{\stackrel{\rightarrow}{#1}}
\newcommand{\toeq}{\stackrel{=}{\to}}
\newcommand{\todesc}{\stackrel{\downarrow}{\to}}
\def\dsp{\displaystyle}
\newcommand{\tosub}[1]{\raisebox{-7.0pt}
        {$\ \stackrel{\textstyle\to}{\scriptstyle#1}\ $}}
\def\rule#1#2{ \dsp \frac{#1}{#2}}

\def\grtypecolor#1{\textcolor{blue}{#1}}
\def\hiordcolor#1{\textcolor{cyan}{#1}}
\def\jf#1#2#3{{#1}, \grtypecolor{#2} |\hiordcolor{#3}}

\def\anyint{?^{\tt int}}
\def\anybool{?^{\tt bool}}
%\def\terminates{\lightning}
\def\tyjf#1#2{\Gamma \vdash #1:#2}

\newtheorem{defn}{Definition}[section]

%%%%%%%%%%%%%%%%%%%%% END OF MY DEFINITIONS

\title{Termination Analysis of a subset of CoreML}
\author[W. Blum]{William Blum\\
\href{mailto://william.blum@comlab.ox.ac.uk}{\tiny william.blum@comlab.ox.ac.uk} }
\institute[University Of Oxford]{Oxford University Computing Laboratory}
\date{BCTCS Nottingham}

\begin{document}


\frame{\titlepage}

\section<presentation>*{Outline}

\begin{frame}
  \frametitle{Outline}
  \tableofcontents[part=1]
\end{frame}

\AtBeginSection[] {
   \begin{frame}<beamer>
     \frametitle{Outline}
     \tableofcontents[currentpart,currentsection]
   \end{frame}
 }


\part<presentation>{Main Talk}

\section[Size-change Principle for first-order programs]{Size-change Principle for first-order programs}

% ==================================================================================

\begin{frame}[fragile]
\frametitle{First order programs}
    \small

 \begin{columns}
   \column{7cm}
\textbf{Untyped functional language} recursion, if-then-else,
primitive operators, single data type
 \vspace{0.4cm}

\uncover<2->{
\textbf{Call-by-value evaluation semantics}: \\
$\begin{array}{ll}
\mathcal{E} \sem{\tt f} \vec{x} = v &  \mbox{{\tt f} evaluates to {\tt v} on input $\vec{x}$}, \\
\mathcal{E} \sem{\tt f} \vec{x} = \bot &  \mbox{{\tt f} does not terminate on input $\vec{x}$}.
  \end{array}
$
\vspace{0.4cm}
}

\uncover<3->{
\textbf{Exact call semantics}: a computation is described by a state transition sequence.
\vspace{0.6cm}
}
\uncover<4->{
\textbf{Finite approximation} of the call semantics: the control flow graph.
\vspace{0.4cm}
}

  \column{4.3cm}
\begin{block}{Example}
\begin{verbatim}
f(x) = gcd(x,18)
gcd(x,y) =
 if y == 0 then x
 else gcd(y, x mod y)
\end{verbatim}

\end{block}

\vspace{0.1cm}

\uncover<3->{
\begin{block}{}
 $f,4 \to gcd,(4,18) \to gcd,(18,2) \to gcd,(2,2) \to gcd,(2,0)$
\end{block}
\vspace{0.1cm}
}

\visible<4->{
\begin{block}{}
\begin{center} \includegraphics{flow0}
\end{center}
\end{block}
}
\end{columns}
\end{frame}

%%%%%%%%%%%%%%%%%%%%%%%%%%%%%%%%%%%%%%%%%%%%%%%%%

\frame{ \frametitle{Termination}

    \begin{block}{Characterization of termination}
    P terminates on all input values \\
$\iff$ Infinite state transition sequences are invalid computations.
    \end{block}
\note[item]{Characterization of termination is proved  by analysis of the definition of the evaluation semantics}

\begin{itemize}
\uncover<2->{
\item \emph{What is an invalid computation?} \ \\
For instance: a computation in which some positive integer variable
decreases infinitely... }

\uncover<3->{
\item \emph{The Size-Change Principle} proves that
for any computation corresponding to an infinite path in the control
flow, the value of some well-founded variable decreases infinitely.
\end{itemize}}
}


%%%%%%%%%%%%%%%%%%%%%%%%%%%%%%%%%%%%%%%%%%%%%%%%%

\frame{ \frametitle{Size-change graphs (SCG)}

\textbf{Definition}: A SCG describes a program call. It consists of
a {\color{red} source set of vertices}, a {\color{blue} target set
of vertices} and a set of labeled arcs.


The SCG
\def\objectstyle{\scriptstyle}
\def\labelstyle{\scriptstyle}
$\left( \vcenter {\xymatrix  @R=0.3cm {
 {\color{red} x} \ar[r]^{=} & {\color{blue} x} \\
*{} & {\color{blue} y } }} \right)$ describes the call from {\tt f}
to {\tt gcd}.


\textbf{Safety:} $\todesc$ arcs denote decreases in parameter value, $\toeq$ arcs denote non increase in parameter value.

\note[item]{
A SCG $G$ \textbf{safely} describes a call if:
    \begin{itemize}
    \item when $x \todesc y \in G$, the value used for the call parameter {\tt y} is strictly less than
    the value of input parameter {\tt x}.

    \item when $x \toeq y \in G$, the value used for the call parameter {\tt y} is
    not greater than the value of input parameter {\tt x}.
    \end{itemize}
}
\vspace{0.2cm}

\emph{Example}: consider the call ``$\tt gcd(\underbrace{\color{red}{\tt y}}_{\color{blue}{\tt x}},
\underbrace{\color{red}{\tt x} {\tt\ mod\ } \color{red}{\tt y}}_{\color{blue}{\tt y}})$'':


$ \left( \vcenter {\xymatrix @R=0.6cm
 {
 {\color{red} x} \ar[rd]^{=} & {\color{blue} x} \\
 {\color{red} y} \ar[r]^{\downarrow} \ar[ur]^{=} & {\color{blue} y }  }} \right)
 \qquad  \left( \vcenter {\xymatrix @R=0.6cm {
 {\color{red} x} \ar[rd]^{=} & {\color{blue} x} \\
 {\color{red} y} \ar[r]^{=} & {\color{blue} y}  }} \right)
 \qquad  \left( \vcenter{ \xymatrix @R=0.6cm {
 {\color{red} x} \ar[r]^{=} & {\color{blue} x} \\
 {\color{red} y} \ar[r]^{\downarrow} & {\color{blue} y}  }} \right)
 \qquad  \left( \vcenter{ \xymatrix @R=0.6cm {
 {\color{red} x}  & {\color{blue} x} \\
 {\color{red} y} & {\color{blue} y} }} \right)
$

Only one of these SCG is not safe for this call. \note[item]{the
scg 1 2 4 are safe for this specific call}

}

%%%%%%%%%%%%%%%%%%%%%%%%%%%%%%%%%%%%%%%%%%%%%%%%%

\frame{ \frametitle{Composition of size-change graphs}

$$ \mbox{If } f \stackrel{G_1}{\longrightarrow} g \mbox{ and } g \stackrel{G_2}{\longrightarrow} h \mbox{ then }
f \stackrel{G_1;G_2}{\longrightarrow} h$$
\begin{eqnarray*}
\overbrace{\xymatrix {
 a \ar[r]^{=} \ar[rd]^{=} & x \ar[rd]^{=} \ar[r]^{\downarrow} & u \\
*{} & y \ar[r]^{\downarrow} & v }}^{G_1 \qquad G_2} \quad & \equiv
& \quad \overbrace{\xymatrix {
 a \ar[r]^{\downarrow} \ar[rd]^{=} &  u \\
*{} & v }}^{G_1 ; G_2}
\end{eqnarray*}

If $\mathcal{G}$ is a set of size-change graphs then $\overline{\mathcal{G}}$ denotes the composition closure of
$\mathcal{G}$.

}

%%%%%%%%%%%%%%%%%%%%%%%%%%%%%%%%%%%%%%%%%%%%%%%%%

\frame{ \frametitle{Size-change termination (SCT)}


\textbf{Definition} Consider $\mathcal{G}$ a set of size-change
graphs.

A program {\tt P} is $\mathcal{G}$-SCT if

\begin{itemize}
  \item $\mathcal{G}$ safely describes {\tt P} (for every reachable call $c$ there is a corresponding SCG
  $G_c \in \mathcal{G}$)
  \item for all infinite computation $cs = \langle c_0 c_1 \ldots \rangle$,
any sequence of size-change graphs $G_{c_0} G_{c_1} \ldots$ (describing safely the calls of $cs$) has an infinite descending thread.
$$\xymatrix @R=0.6cm {
 x & u \ar[rd]^{=} & u & x \ar[r]^{=} & \ldots\\
 y & v  & v \ar[ru]^{\downarrow} & y & \ldots
}
$$
\end{itemize}

\note[item] {P is G-SCT means that G is fine enough to describe descents in infinite computation.}
\note[item] {A thread is a sequence of arcs starting from any graph and going through the consecutive graphs in
the sequence.}
\note[item] {Descending means that there is at least one arc in the thread of type $\downarrow$}


\uncover<2->{
We assume that data-types are well-founded.
\note[item]{well-founded = every non-empty subset E of X has a minimal element.}

\begin{block}{Theorem}
If {\tt P} is $\mathcal{G}$-SCT then {\tt P} terminates for all input values
\end{block}
}
\note[item]{The theorem is a consequence of the definition: all infinite path in the CFD are impossible computation therefore
P is terminating.}

}

\frame{ \frametitle{Deciding Size-Change Termination}
%%%%%%%%%%%%%%%%%%%%%%%%%%%%%%%%%%%%%%%%%%%%%%%%%


\begin{block}{$\mathcal{G}$-SCT characterization [Jones et al. 2001]}
$${\tt P} \mbox{ is \textbf{not} $\mathcal{G}$-SCT} $$
$$\iff$$
$$ \exists {\tt f} \stackrel{G}{\to} {\tt f} \in \overline{\mathcal{G}}
        \mbox{ such that }
 \left( \begin{array}{l} G;G = G \\ \forall {\tt x} \in gb({\tt f}) : {\tt x} \todesc {\tt x} \not \in G \end{array} \right)
$$
\end{block}

Hence $\mathcal{G}$-SCT is decidable. And it is PSPACE-complete (see \cite{jones01})

  \note[item]{The number of composition can be exponential in the size of the program size.}
  \note[item]{Algorithm for deciding $\mathcal{G}$-SCT:
    \begin{itemize}
      \item describe safely each call by a SCG. This gives a set $\mathcal{G}$ of scgs.
      \item compute the composition closure of $\mathcal{G}$
      \item check the presence of non-descending asymptotic loop size-change graph
    \end{itemize}
  }
}

%%%%%%%%%%%%%%%%%%%%%%%%%%%%%%%%%%%%%%%%%%%%%%%%%

\section{An extension for a subset of Core ML}


\frame{ \frametitle{The language $\mathcal{L}_{ml}$}

\textbf{Grammar:}
{\small
\begin{center}
\begin{tabular}{ll}
{\tt e ::= x, f}     & value identifiers \\
{\tt\ \ \ \ | true | false } & boolean constants \\
{\tt\ \ \ \ | if e then e else e} & conditional \\
{\tt\ \ \ \ | n} & integer constants (n $\in \natset$)\\
{\tt\ \ \ \ | e = e} & integer equality \\
{\tt\ \ \ \ | succ e | pred e } & successor and predecessor\\
{\tt\ \ \ \ | \alert{fun (x:ty) -> e}} & function abstraction \\
{\tt\ \ \ \ | \alert{fun f=(x:ty) -> e}} & recursively defined function \\
{\tt\ \ \ \ | e e } & function application \\
{\tt\ \ \ \ | let x = e in e } & local variable definition
\end{tabular}
\end{center}
}

A program is a \alert{single} closed expression. \note[\item]{ {\tt
P} is a closed expression, so how to prove termination for any input
value? We will see that it is not a problem in our framework,
fortunately.}

\vspace{0.4cm} \textbf{Data types:} ground values +
\alert{higher-order functions}.

\note[\item]{note that we restrict the datatype to positive integers in order to have a well-founded set.}
}

%%%%%%%%%%%%%%%%%%%%%%%%%%%%%%%%%%%%%%%%%%%%%%%%%


\frame{ \frametitle{Semantics of $\mathcal{L}_{ml}$ (environment
based)}

\emph{Canonical expressions}: $\natset \union \boolset \union \{ {\tt e} \st {\tt e} \mbox{
is an abstraction} \}$
\begin{eqnarray*}
State &=& \{ {\tt e}:\rho \st {\tt e} \in subexp({\tt P}), \rho \in Env, fv({\tt e}) \subseteq dom(\rho) \} \\
Value &=& \{ {\tt e}:\rho \in State \st \mbox{{\tt e} canonical} \}\\
Env &=& \{ \rho : X \to Value \st X \mbox{ finite set of variables} \}
\end{eqnarray*}

Let $s \in State$, $v \in Value$ and $\rho \in Env$

\textbf{Call-by-value evaluation semantics} ``$s\ \evaluates\ v$''
$$\tt\rule{}{v:\rho\ \evaluates\ v:\rho} (v \mbox{ canonical})$$

\textbf{Run-time errors} ``$s\ \error$'' \qquad
$\mbox{(ErrOp1)}\tt \rule{e:\rho\ \evaluates\ 0}{pred\ e:\rho\
\error}$

\note[item]{``$s\ \error$'' means ``trying to evaluate $s$ will
raise an error and therefore the computation terminates.''}
\note[item]{Error handling is important: this permits to detect the
program \texttt{loop (desc x)} as terminating. Where \texttt{loop}
is a program looping infinitely and \texttt{desc} is defined as
``desc x=desc(x-1)'' (which is terminating since it raises an
error).}

\textbf{Call semantics} ``$ s \textcolor{red}{\to} s'$''
$$\mbox{(CallG)} \rule{ \tt e_1:\rho\ \evaluates\ \mbox{\tt fun (x:ty)->} e_0 : \rho_0 \qquad
                \tt e_2 : \rho\ \evaluates\ v_2
                }
                {
                \tt e_1 e_2 : \rho \textcolor{red}{\tosub{c}} e_0 : \rho_0 [x\mapsto v_2]
                }
$$

\note[item]{
What is a call? Where do the calls occur?

`` e1 $\to$ e2 if it is necessary to know an evaluation of e2 in order to deduce
an evaluation of e1'', \cite{jones04}
}
}


\frame{ \frametitle{Graph generation}

\begin{itemize}
\item Two SCG generated per call:
$\hiordcolor{G^+}$ describing higher-order values and
$\grtypecolor{G^0}$ for ground type values. \vspace{0.2cm}

\item The \emph{free variables} of an expression correspond to the \emph{input parameters} in the first-order case.
\vspace{0.2cm}

\item We define well-founded notions of size for higher-order and ground type expressions.
\vspace{0.2cm}

\item We extend the semantic rules to generate safe SCG:
$$\mbox{(ValueG)} \rule{}{
        \jf{\tt v\ \evaluates\ v}{id_{\tt e}^{=}}{id_{\tt e}^{=}}
        }
        ({\tt v = e:\rho} \mbox{ in canonical form})
$$

\begin{bigcenter}
(CallG)
$\rule{ \jf{\tt e_1:\rho\ \evaluates\ \mbox{\tt fun
(x:ty)->} e_0 : \rho_0} {G_1}{G_1^+} \qquad
                \jf{\tt e_2 : \rho\ \evaluates\ v_2}{G_2}{G_2^+}
                }
                {
                \jf{\tt e_1 e_2 : \rho \textcolor{red}{\tosub{c}}
         e_0 : \rho_0 [x\mapsto v_2]}
                {CallGr^0_x(G_1,G_2)} {CallGr^+_x(G_1^+,G_2^+)}
                }
$
\end{bigcenter}

\end{itemize}
}



%%%%%%%%%%%%%%%%
\iflongonly \frame{ \frametitle{Safety of SCG for higher-order
values}
\begin{defn}[Size relation $\succeq$]
\label{def:lambda_size}
\begin{center}
$\overbrace{e_1 : \rho_1}^{s_1} \succeq \overbrace{e_2 : \rho_2}^{s_2}$ \\
$\iff$ \\
$  support(s_1) \ni s_2
  \zor \left( subexp(e_1) \ni e_2\ \zand\ {\rho_1}_{|fv(e_2)} = \rho_2 \right)
$
\end{center}

\end{defn}
where: \\
$$ support({\tt e}:\rho) = \{ {\tt e} : \rho \} \union \bigcup_{{\tt x} \in fv({\tt e})} support(\rho({\tt x}))$$
}

\frame{ \frametitle{Safety of SCG for ground-type values}
\note[item]{$\geq_{\tt int}$ denotes a well-founded order on the integers
and $\succeq_{\tt bool}$ denotes a well-founded order on the boolean $\{ true, false\}$.}
\begin{defn}[Size relation $\succeq_0$]\label{def:ml0_size}
\begin{center}
$\overbrace{e_1 : \rho_1}^{s_1} \succeq_0
\overbrace{e_2 : \rho_2}^{s_2}$
\end{center}
$$\triangleq$$
$$\begin{array}{ll}
 \exists\ \Gamma \mbox{ s.t. } &
             \left\{ \begin {array}{lll}
             & \tyjf{e_1}{int}, \qquad \tyjf{e_2}{int}\\
            & \mbox{ and }
            s_1 \Downarrow n_1  \zand s_2 \Downarrow n_2
            \implies n_1 \geq_{\tt int} n_2
            \end{array} \right. \\
            & or \\
            &
             \left\{ \begin {array}{lll}
                & \tyjf{e_1}{bool}, \qquad \tyjf{e_2}{bool}\\
                \mbox{ and } &
                s_1 \Downarrow b_1  \zand s_2 \Downarrow b_2
            \implies b_1 \succeq_{\tt bool} b_2
            \end{array} \right. \\
            & or \\
            & \tyjf{e_1}{ty_1 \to ty_2} \quad\mbox{ and }\quad \tyjf{e_2}{ty_3 \to ty_4}
\end{array}
$$
\end{defn}
}
\else
\fi
%%%%%%%%%%%%%%%%%%%%%%%%%%%%%


\frame{ \frametitle{Finite approximation of the call semantics}

\textbf{We need a ``control flow graph'' for ML programs} \\

Solution:
  \begin{itemize}
    \item drop the $\rho$ components of the states
    \item abstract integers by a single symbol ``$\anyint$''.
  \end{itemize}
  We obtain a finite abstraction of the computation.

  The set of vertices of the control flow graph (i.e control points) is:
   $$\mathcal{P} = subexp({\tt P}) \union \{ \anyint \}$$

}

%%%%%%%%%%%%%%%%%%%%%%%%%%%%%
\iflongonly
\frame{ \frametitle{Approximation (2)}

\begin{lem}[Approximation]
\label{lem:ml_approx} Suppose ${\tt P}:[] \to^* {\tt e}:\rho$ then
for ${\tt e'} \in subexp({\tt P})$, $n \in \natset$, $b \in \boolset $ and $v$ an abstraction:
\vspace{0.3cm}

\begin{tabular}{lll}
Exact semantics & & Approximation semantics\\
$\jf{\tt e:\rho \to e':\rho'}{G^0}{G^+}$ & $\implies$ & $\jf{\tt e \to e'}{G^0}{G^+}$
\\
$ \jf{{\tt e}:\rho\ \evaluates\ n:[]}{G^0}{G^+}$ &
$\implies$& $\jf{\tt e\ \evaluates\ \anyint}{G^0}{G^+}$
\\
 $\jf{{\tt e}:\rho\ \evaluates\ b:[]}{G^0}{G^+}$  & $\implies$& $\jf{\tt e\ \evaluates\
\anybool}{G^0}{G^+}$
\\
 $ \jf{{\tt e}:\rho\ \evaluates\ v:\rho'}{G^0}{G^+}$  & $\implies$& $\jf{\tt e\ \evaluates\
v}{G^0}{G^+}$
\end{tabular}
\end{lem}
}
\else
\fi
%%%%%%%%%%%%%%%%%%%%%%%%%%%%%

\frame{ \frametitle{The size-change principle}

What do we have:
\begin{itemize}
\item Termination characterized by infinite call sequences
\item Well-founded order on the data values
\item Finite approximation of call semantics
\item We can compute two safe sets of size-change graphs describing the calls
(by applying the semantic rules exhaustively).
\end{itemize}

\note[item]{$\mathcal{G}$ can be computed by applying exhaustively
the semantic rules starting with expression P until no new graphs or
subexpressions are obtained. This process ends since P contains a
finite number of subexpressions and possible size-change graphs.}

\textbf{Hence the SCP can be applied! (twice)}

}

\frame{ \frametitle{Results}

\begin{itemize}
\item Counter example

\note[item]{counter-example page 65 of the thesis}
\begin{semiverbatim}
let rec counter x =

\ if x = 0 then counter (succ x) else 1

in counter 7;;
\end{semiverbatim}
is terminating but not SCT.

\item Ackerman's function: SCT relatively to ground-type values.

\item Function computing the minimum of two numbers:
    \begin{itemize}
    \item is \alert{SCT} if we use the native representation of
    integers provided by $\mathcal{L}_{ml}$,

    \item \alert{is not SCT} if we use Church numeral to encode integers.
    \end{itemize}

\end{itemize}
}


\section*{Summary}

\frame{ \frametitle{Conclusion}

\begin{itemize}
\item The Size-Change Principle from Neil D. Jones et al.
    \begin{itemize}
      \item based on a finite approximation of the call semantics,
      \item and a safe description of the calls.
    \end{itemize}
\item Extension to a higher-order functional language
    \begin{itemize}
    \item detects decrease on ground-type values as well as higher-order values
    \item allows local definition {\tt let}
    \item handles recursion natively (no need to define a Y combinator)
    \item handles numbers natively
    \end{itemize}
\item Further direction

sequential composition,
storage location and references,
tuples
list,
user defined structures,
for and while loop structures.
\end{itemize}
}

\subsection<presentation>*{Bibliography}

\begin{frame}
  \frametitle<presentation>{Bibliography}

  \begin{thebibliography}{10}

  \beamertemplatearticlebibitems
  \bibitem{jones01}
    Chin~Soon~Lee, Neil~D.~Jones, and Amir~Ben-Amram.
    \newblock The Size-Change Principle for Program Termination.
    \newblock {\em Principles of Programming Languages}, pp. 81-92. Volume 28 of Principles of Programming Languages. ACM press 2001

  \bibitem{jones04}
    Neil~D.~Jones and Nina~Bohr
    \newblock Termination Analysis of the Untyped
    $\lambda$-Calculus.
    \newblock {\em Rewriting Techniques and Applications.},Proceedings (V van Oostrom, ed.), pp. 1-23. Volume 3091 of LNCS. Springer-Verlag 2004.


  \bibitem{pitts2002}
    Andrew~M.~Pitts
    \newblock Operational Semantics and Program Equivalence.
    \newblock {\em Applied Semantics.}, pp. 378-412. Volume 2395 of LNCS. Springer-Verlag 2002

  \end{thebibliography}
\end{frame}

\end{document}
